\chapter{Local Consistency Residual and Jacobian as Independent Customization Points}
\label{ch:local-consistency-residual-and-jacobian}

\section{Motivation}

The previous material refactors had already separated:
\begin{itemize}
  \item the \emph{yield function} from the yield-surface geometry;
  \item the \emph{return algorithm} from the constitutive ingredients;
  \item the \emph{constitutive state} from the low-level storage backend.
\end{itemize}

However, one coupling still remained inside the local plastic solve: the
algorithm knew too much about the specific scalar equation being solved. In the
classical radial-return setting this is easy to hide because one often writes a
closed-form update for the plastic multiplier increment $\Delta\gamma$. That
shortcut does not scale well to future models.

In particular, the next generations of models needed by this library:
\begin{itemize}
  \item pressure-sensitive plasticity,
  \item nonlinear hardening laws,
  \item finite-strain multiplicative plasticity,
  \item damage-plasticity and crack-regularised constitutive updates,
\end{itemize}
all tend to be expressed as local nonlinear residual/Jacobian problems rather
than as one hard-coded formula for $\Delta\gamma$.

\section{Classical Local Plasticity Problem}

In the present small-strain rate-independent setting, the constitutive update is
still organised around the elastic predictor and the plastic corrector:
\begin{align}
  \boldsymbol{\sigma}^{\mathrm{tr}} &= \mathbb{C}_e :
    \left(\boldsymbol{\varepsilon}_{n+1}
          - \boldsymbol{\varepsilon}_n^p\right), \\
  f^{\mathrm{tr}} &= q(\boldsymbol{\sigma}^{\mathrm{tr}})
    - \sigma_y(\alpha_n).
\end{align}

If $f^{\mathrm{tr}} \le 0$, the step is elastic. Otherwise the constitutive
site must solve a local nonlinear problem for the plastic correction. In the
simplest radial-return case this can be written as a scalar consistency
equation
\[
  r(\Delta\gamma) = 0.
\]

The important architectural observation is that this scalar equation is not the
same concept as the nonlinear algorithm used to solve it.

\section{Architectural Split}

The implementation now distinguishes four independent compile-time roles:
\begin{enumerate}
  \item \textbf{Yield criterion}: geometry of the admissible surface.
  \item \textbf{Yield function}: scalar overstress map.
  \item \textbf{Consistency residual / Jacobian}: definition of the local
        nonlinear equation and its derivative.
  \item \textbf{Return algorithm}: the nonlinear solve strategy itself.
\end{enumerate}

This leads to the following staging:
\[
  \texttt{YieldCriterion}
  \;\to\;
  \texttt{YieldFunction}
  \;\to\;
  \texttt{ConsistencyResidual/Jacobian}
  \;\to\;
  \texttt{ReturnAlgorithm}.
\]

That separation is more than naming hygiene. It is exactly the pattern needed
when the library later moves from scalar J2 updates to local residual/Jacobian
systems in finite strains.

\section{Default Residual}

The new default residual evaluates the corrected state semantically rather than
re-deriving the current closed-form denominator by hand. It computes:
\begin{enumerate}
  \item corrected stress
    $\boldsymbol{\sigma}(\Delta\gamma)$;
  \item evolved internal variables
    $\alpha(\Delta\gamma)$;
  \item corrected trial state reconstructed from
    $\boldsymbol{\sigma}(\Delta\gamma)$;
  \item effective state after optional backstress shift;
  \item residual
    \[
      r(\Delta\gamma)
      = f\!\left(\boldsymbol{\sigma}(\Delta\gamma), \alpha(\Delta\gamma)\right).
    \]
\end{enumerate}

This is slower than a hand-derived one-line formula, but it is the right
default seam because it keeps the residual faithful to the constitutive model
rather than to one particular radial-return algebra.

\section{Default Jacobians}

Two Jacobian paths are now available.

\subsection*{Closed-form Jacobian}

The fast default for the present associated J2 family remains a closed-form
derivative:
\[
  r'(\Delta\gamma)
  = -\sqrt{\frac{2}{3}}\left(3G + H\right),
\]
evaluated through the current constitutive kernel.

\subsection*{Finite-difference Jacobian}

A generic numerical derivative is also available:
\[
  r'(\Delta\gamma)
  \approx
  \frac{r(\Delta\gamma + h) - r(\Delta\gamma - h)}{2h}.
\]

This is not the preferred production path for hot local loops, but it is
important for development:
\begin{itemize}
  \item it provides a correctness baseline;
  \item it lets a new material model become usable before its exact local
        Jacobian has been derived;
  \item it mirrors the type of prototyping stage often needed for complex
        finite-strain or damage models.
\end{itemize}

\section{Newton Consistency Algorithm}

With residual and Jacobian split out, the library now provides a generic Newton
solver for the scalar consistency problem:
\begin{align}
  r(\Delta\gamma_k) &= 0, \\
  \Delta\gamma_{k+1}
    &= \Delta\gamma_k - \frac{r(\Delta\gamma_k)}{r'(\Delta\gamma_k)}.
\end{align}

The initial guess reuses the current closed-form estimate of the J2 radial
return. This is intentional:
\begin{itemize}
  \item for the present model family it gives a high-quality starting point;
  \item for future model families it still keeps the Newton solver as a clean,
        reusable algorithmic shell.
\end{itemize}

\section{Why This Matters for Finite Strains}

Finite-strain plasticity is rarely implemented as a direct analogue of the
small-strain radial return. Even when the unknown remains scalar, the local
problem is usually expressed in terms of a residual/Jacobian pair attached to:
\begin{itemize}
  \item multiplicative splits,
  \item Kirchhoff or Mandel stresses,
  \item logarithmic or objective rates,
  \item internal tensorial variables.
\end{itemize}

From that perspective, the present refactor is strategically important:
\begin{quote}
the code now has an explicit place where the local nonlinear constitutive
equation lives, and that place is independent from the nonlinear solver.
\end{quote}

That is the exact abstraction needed to grow from the current small-strain J2
path into finite-strain plasticity and eventually damage/fracture regularized
models.

\section{Verification}

The regression suite now checks three things:
\begin{itemize}
  \item a Newton-based local solve reproduces the current standard radial
        return for the existing J2 family;
  \item the same Newton algorithm remains usable when the Jacobian is replaced
        by a finite-difference policy;
  \item changing only the consistency residual changes the final constitutive
        response, proving that the residual is now a real customization point.
\end{itemize}

\section{Architectural Outcome}

After this step, the local constitutive update is no longer just:
\[
  \texttt{YieldFunction} + \texttt{ReturnAlgorithm}.
\]

It is now:
\[
  \texttt{YieldFunction}
  + \texttt{ConsistencyResidual}
  + \texttt{ConsistencyJacobian}
  + \texttt{ReturnAlgorithm}.
\]

That is a much better foundation for the next stages:
\begin{itemize}
  \item finite-strain plasticity with explicit continuum state;
  \item local residual/Jacobian systems of dimension larger than one;
  \item pressure-sensitive and damage-coupled return algorithms;
  \item fracture-oriented local solvers and regularization interfaces.
\end{itemize}

\section{References}

\begin{thebibliography}{9}

\bibitem{simo1998local}
J.~C. Simo and T.~J.~R. Hughes,
\emph{Computational Inelasticity},
Springer, 1998.

\bibitem{desouza2008}
E.~A. de Souza Neto, D.~Peri\'{c}, and D.~R.~J. Owen,
\emph{Computational Methods for Plasticity: Theory and Applications},
Wiley, 2008.

\bibitem{crisfield1997}
M.~A. Crisfield,
\emph{Non-Linear Finite Element Analysis of Solids and Structures, Vol.~2:
Advanced Topics},
Wiley, 1997.

\end{thebibliography}
